\SGAchapter{XI}{完全交的上同调}{Cohomologie des intersections complètes}
\begin{center}P. Deligne\end{center}
\SourcePage{49}
\SGAsection{1}{射影空间与完全交的上同调}
设 $X$ 是纯 $n$ 维的复\Term{nonsingular}{非奇异}\Term{projective-variety}{射影簇}，$\mathcal O(1)$ 是 $X$ 上的一个\Term{ample}{丰沛}\Term{invertible-sheaf}{可逆层}。根据 Akizuki 与 Nakano [1]，由 Kodaira 的\Term{vanishing-theorem}{消没定理} [4] 可得
\[
 H^i(X,\Omega_X^j(-1))=0\qquad\text{当 }i+j<n.\tag{*}
\]
仍根据 [1]，这一结果蕴含复\Term{cohomology}{上同调}中的“\Term{weak-lefschetz-theorem}{弱 Lefschetz 定理}”。我们拟用直接计算证明：若 $X$ 是任意\Term{characteristic-field}{特征}的域上的一个\Term{projective-space}{射影空间}，则 (*) 仍成立。这样便能仿照 [1]，把任意\Term{characteristic-field}{特征}的域上\Term{complete-intersection}{完全交}的\Term{hodge-cohomology}{Hodge 上同调}和\Term{de-rham-cohomology}{De Rham 上同调}与\Term{projective-space}{射影空间}的相应\Term{cohomology}{上同调}作比较。

随后，我们将具体说明所得结果的若干推论。在第 2 节中，给出由 Hirzebruch 的数值公式 (2.2) \Term{derive}{导出}的定性结果。我要感谢 N. Katz 在这一节的构思中给予我的帮助。

\textbf{1.0。}设 $E$ 是\Term{scheme}{概形} $S$ 上\Term{rank}{秩}为 $r+1$ 的\Term{locally-free}{局部自由}模。假设 $r>0$。回忆 $\mathbf P(E)$ 表示 $S$ 上参数化 $E$ 的\Term{hyperplane}{超平面}的\Term{relative-projective-space}{相对射影空间}。以 $p$ 表示 $\mathbf P(E)$ 到 $S$ 的\Term{projection}{投影}，以
\[
 \eta\in H^0(S,R^1p_*\Omega^1_{\mathbf P(E)/S})
\]
表示\Term{invertible-sheaf}{可逆层} $\mathcal O(1)$ 的 Hodge 意义的\Term{first-chern-class}{第一 Chern 类}。
\SourcePage{50}
回忆，$\eta$ 是 $\mathcal O(1)$ 在 $H^0(S,R^1p_*\mathcal O^*)$ 中的类在\Term{logarithmic-differential}{对数微分}
\[
 df/f:\mathcal O^*\longrightarrow\Omega^1_{\mathbf P(E)/S}
\]
下的像。对于 $\mathbf P(E)$ 上任意 $\mathcal O$ 模\Term{sheaf}{层} $F$，令
\[
 F(n)=F\otimes\mathcal O(1)^{\otimes n}\qquad(n\in\mathbf Z).
\]
\begin{SGAresult}{\Term{theorem}{定理}}{1.1}
(i) \Term{sheaf}{层} $R^ip_*(\Omega^j_{\mathbf P(E)/S}(n))$ 是\Term{locally-free}{局部自由}的，其形成与任意\Term{base-change}{基变换}相容。

(ii) 对于 $0\leq i\leq r$，$R^ip_*(\Omega^i_{\mathbf P(E)/S})$ 的\Term{section}{截面} $\eta^i$ 给出此\Term{sheaf}{层}与 $\mathcal O_S$ 之间的\Term{isomorphism}{同构}。此外，
\[
 R^ip_*\Omega^j_{\mathbf P(E)/S}=0\qquad\text{当 }i\ne j\text{ 或 }i>r.
\]
(iii) 若 $n\ne0$，则除下列两种情形外，$R^ip_*(\Omega^j_{\mathbf P(E)/S}(n))$ 为零：
\[
 \begin{array}{ll}
 (\mathrm{iii}_a)&i=0\text{ 且 }n>j,\\
 (\mathrm{iii}_b)&i=r\text{ 且 }n<j-r.
 \end{array}
\]
\end{SGAresult}
\begin{proof}
设 $\mathbf V(E)=\operatorname{Spec}(\operatorname{Sym}^*E)$ 是由 $E$ 定义的\Term{affine-space}{仿射空间}，$e$ 为其\Term{zero-section}{零截面}，$\mathbf V(E)^*=\mathbf V(E)-e(S)$ 为其补集，$a$ 与 $a'$ 分别为 $\mathbf V(E)$ 和 $\mathbf V(E)^*$ 到 $S$ 的\Term{projection}{投影}，$q$ 为 $\mathbf V(E)^*$ 到 $\mathbf P(E)$ 的\Term{projection}{投影}：
\[
\begin{tikzcd}
\mathbf V(E)\arrow[dr,"a"'] & \mathbf V(E)^*\arrow[l,hook']\arrow[r,"q"]\arrow[d,"a'" description] & \mathbf P(E)\arrow[dl,"p"]\\
& S\arrow[ul,"e",bend left=35] &
\end{tikzcd}
\]
$\mathbf P(E)$ 上去掉\Term{zero-section}{零截面}的\Term{affine-space}{仿射空间} $\mathbf V(\mathcal O(1))^*$ \Term{canonical}{典范}地\Term{isomorphism}{同构}于 $\mathbf V(E)^*$，从而对 $\mathbf P(E)$ 上任意\Term{quasi-coherent}{拟凝聚} $\mathcal O$ 模\Term{sheaf}{层} $F$，\Term{canonical}{典范}地有
\SourcePage{51}
\[
 q_*q^*F\simeq\bigoplus_{n\in\mathbf Z}F(n).\tag{1.1.1}
\]
$q_*q^*F$ 的这个\Term{grading}{分次}可以从 $\mathbf V(E)$ 的\Term{homothety}{位似}群 $\mathbf G_m$ 在此\Term{sheaf}{层}上的\Term{action}{作用}\Term{derive}{导出}。

$\mathbf V(E)$ 上的\Term{relative-vector-field}{相对向量场}可与从 $\mathbf V(E)$ 到 $\mathbf V(E)$（后者视为原点的\Term{tangent-space}{切空间}）的 $S$-\Term{morphism}{态射}相等同。由 $\mathbf V(E)$ 的\Term{identity-map}{恒等映射}定义的\Term{vector-field}{向量场} $X$ 在\Term{homothety}{位似}下\Term{invariant}{不变}。此\Term{vector-field}{向量场}与 $q$ 的\Term{fibre}{纤维}\Term{tangent}{相切}，并定义一个\Term{trivialization}{平凡化}
\[
 X:\Omega^1_{\mathbf V(E)^*/\mathbf P(E)}\longrightarrow\mathcal O_{\mathbf V(E)^*}.\tag{1.1.2}
\]
若 $x_i$ 构成 $E$ 的一组基，则
\[
 X=\sum_i x_i\,dx^i.
\]
由于 $\Omega^1_{\mathbf V(E)^*/\mathbf P(E)}$ 的\Term{rank}{秩}为一，\Term{exact-sequence}{正合列}
\[
0\longrightarrow q^*\Omega^1_{\mathbf P(E)/S}\longrightarrow\Omega^1_{\mathbf V(E)^*/S}\longrightarrow\Omega^1_{\mathbf V(E)^*/\mathbf P(E)}\longrightarrow0
\]
诱导\Term{exact-sequence}{正合列}
\[
0\longrightarrow q^*\Omega^i_{\mathbf P(E)/S}\longrightarrow\Omega^i_{\mathbf V(E)/S}\longrightarrow\Omega^1_{\mathbf V(E)^*/\mathbf P(E)}\otimes q^*\Omega^{i-1}_{\mathbf P(E)/S}\longrightarrow0.
\]
利用 (1.1.2)，可将其写成
\[
0\longrightarrow q^*\Omega^i_{\mathbf P(E)/S}\longrightarrow\Omega^i_{\mathbf V(E)^*/S}\xrightarrow{X\mathbin{\lrcorner}}q^*\Omega^{i-1}_{\mathbf P(E)/S}\longrightarrow0.\tag{1.1.3}
\]
\Term{canonical}{典范}地有
\[
 \Omega^i_{\mathbf V(E)/S}=a^*\bigwedge^iE.\tag{1.1.4}
\]
因此，\Term{short-exact-sequence}{短正合列} (1.1.3) 是将 $\mathbf V(E)^*$ 上的\Term{long-exact-sequence}{长正合列}
\[
 \cdots\longrightarrow a'^*\bigwedge^iE\xrightarrow{X\mathbin{\lrcorner}}a'^*\bigwedge^{i-1}E\xrightarrow{X\mathbin{\lrcorner}}a'^*\bigwedge^{i-2}E\longrightarrow\cdots\tag{1.1.5}
\]
分段而得的。
\SourcePage{52}
此\Term{long-exact-sequence}{长正合列}正是 $\mathbf V(E)$ 上 $\mathcal O_{e(S)}$ 的\Term{koszul-resolution}{Koszul 分解}在 $\mathbf V(E)^*$ 上的\Term{restriction}{限制}。

经由\Term{affine-morphism}{仿射态射} $q$，将\Term{long-exact-sequence}{长正合列} (1.1.5) \Term{projection}{投影}到 $\mathbf P(E)$ 上。由 (1.1.1) 及公式 $q_*a'^*=q_*q^*p^*$，得到\Term{long-exact-sequence}{长正合列}
\[
 \cdots\longrightarrow\bigoplus_{n\in\mathbf Z}p^*\bigwedge^iE(n)\xrightarrow{X\mathbin{\lrcorner}}\bigoplus_{n\in\mathbf Z}p^*\bigwedge^{i-1}E(n)\longrightarrow\cdots.\tag{1.1.6}
\]
它\Term{separable}{可分}段为 (1.1.3) 的像，即\Term{short-exact-sequence}{短正合列}
\[
0\longrightarrow\bigoplus_{n\in\mathbf Z}\Omega^i_{\mathbf P(E)/S}(n)\longrightarrow\bigoplus_{n\in\mathbf Z}p^*\bigwedge^iE(n)\xrightarrow{X\mathbin{\lrcorner}}\bigoplus_{n\in\mathbf Z}\Omega^{i-1}_{\mathbf P(E)/S}(n)\longrightarrow0.\tag{1.1.7}
\]
\Term{exact-sequence}{正合列} (1.1.7) 是下列\Term{exact-sequence}{正合列}之和：
\[
0\longrightarrow\Omega^i_{\mathbf P(E)/S}(n)\longrightarrow p^*\bigwedge^iE(n-i)\xrightarrow{X\mathbin{\lrcorner}}\Omega^{i-1}_{\mathbf P(E)/S}(n)\longrightarrow0.\tag{1.1.8}
\]
回忆（EGA III 2.1.15 或 FAC），$R^ip_*\mathcal O(k)$ 是\Term{locally-free}{局部自由}的，其形成与任意\Term{base-change}{基变换}相容，且
\[
\begin{split}
R^ip_*\mathcal O(k)=0\quad&\text{若 }i\ne0,r,\text{ 或 }i=0,\ k<0,\\
&\text{或 }i=r,\ k>-r-1,
\end{split}\tag{1.1.9}
\]
因此 $R^jp_*p^*\bigwedge^iE(k)$ 在同样的条件下也为零。

\Term{exact-sequence}{正合列} (1.1.8) 给出 $\Omega^i_{\mathbf P(E)/S}(n)$ 的一个\Term{resolution}{分解}
\[
0\longrightarrow\Omega^i_{\mathbf P(E)/S}(n)\xrightarrow{E}p^*\bigwedge^iE(n-i)\xrightarrow{X\mathbin{\lrcorner}}\cdots\longrightarrow p^*\bigwedge^0E(n-0)\longrightarrow0.\tag{1.1.10}
\]
由\Term{degree}{次数}的理由 (1.1.9)，此\Term{resolution}{分解}定义的\Term{spectral-sequence}{谱序列}\Term{degeneration}{退化}（$E_1=E_\infty$）；对于 $0<j<r$，它给出
\[
R^jp_*\Omega^i_{\mathbf P(E)/S}(n)=\mathscr H\bigl(\cdots\xrightarrow{X\mathbin{\lrcorner}}p_*\bigwedge^{i-j}E(n-i+j)\xrightarrow{X\mathbin{\lrcorner}}\cdots\bigr).\tag{1.1.11}
\]
由于 $r\geq1$，由一个\Term{depth}{深度}论证可知
\[
\begin{split}
\bigoplus_{n\in\mathbf Z}p_*p^*\bigwedge^iE(n)
&=p_*q_*q^*p^*\bigwedge^iE\\
&=a'_*a'^*\bigwedge^iE=a_*a^*\bigwedge^iE.
\end{split}\tag{1.1.12}
\]
\SourcePage{53}
所以 $p_*p^*\bigwedge^iE(n)$ 是 $a_*a^*\bigwedge^iE$ 中权为 $n+i$ 的\Term{homogeneous-component}{齐次部分}，即 $\operatorname{Sym}^n(E)\otimes\bigwedge^iE$。回忆 (1.1.5) 来自\Term{koszul-resolution}{Koszul 分解}，由 (1.1.11) 可知：对于 $0<j<r$，$R^jp_*\Omega^i_{\mathbf P(E)/S}(n)$ 仅在 $i-j=0$ 且 $n-i+j=0$，即 $n=0$ 且 $i=j$ 时才可能非零。此外，在这种情形下，$R^ip_*\Omega^i_{\mathbf P(E)/S}$ \Term{isomorphism}{同构}于 $\mathcal O_S$。

对于 $j=0$，由 (1.1.10) 还可得
\[
p_*\Omega^i_{\mathbf P(E)/S}(n)=\operatorname{Ker}\left(X\mathbin{\lrcorner}:p_*\bigwedge^iE(n-i)\longrightarrow p_*\bigwedge^{i-1}E(n-i+1)\right),\tag{1.1.13}
\]
而同一论证表明，当 $n-i\leq0$ 时，$X\mathbin{\lrcorner}$ 是\Term{injectivity}{单射}。

在 (1.1.8) 中取 $n=0$、$i=r+1$，可得
\[
\Omega^r_{\mathbf P(E)/S}\simeq p^*\bigwedge^{r+1}E(-r-1).
\]
当 $S$ 是域的\Term{spectrum}{谱}时，1.1 (ii)、(iii) 中关于\Term{cohomology-group}{上同调群}为零的断言，在 $j\ne r$ 时由上述论证得到，在 $j=r$ 时则由\Term{duality}{对偶性}（[2] V 11.2，第 211 页）推得。特别地，对任意 $j,n$，至多只有一个 $i$ 使 $R^ip_*\Omega^j_{\mathbf P(E)/S}(n)\ne0$；由于 $\Omega^j_{\mathbf P(E)/S}$ 在 $S$ 上\Term{flat}{平坦}，这就蕴含 (i)（最难的是在 EGA III 中找到引文；我建议：6.10.5 + 7.4.1 + 7.5.5 + 一次取\Term{limit}{极限}）。

已知在 $0\leq i\leq r$ 时 $R^ip_*\Omega^j_{\mathbf P(E)/S}$ 是\Term{invertible-sheaf}{可逆层}，只需证明当 $S$ 为域的\Term{spectrum}{谱}时，$\eta^i\ne0$。事实上，$\eta^r$ 是一个\Term{rational-point}{有理点}（$r$ 个\Term{hyperplane}{超平面}的交）的\Term{cohomology-class}{上同调类}，因此是 $H^r(\mathbf P^r,\Omega^r)$ 的一组基（FGA，第 149 讲，第 4 节）。
\end{proof}

\textbf{1.2。}可以直接从 (1.1.10) 或 (1.1.13) 计算 $\Omega^i(n)$ 的\Term{euler-poincare-characteristic}{Euler--Poincaré 示性数}。所得结果并不十分诱人：
\[
\chi(\mathbf P_k^r,\Omega^j_{\mathbf P_k^r}(n))=
\sum_{a\leq i\leq j}(-1)^{i-j}\binom{r+1}{i}\binom{r+n-i}{r}.
\]

\SourcePage{54}
\begin{SGAresult}{命题}{1.3}
设 $f:X\longrightarrow S$ 是\Term{proper}{固有}\Term{smooth}{光滑}\Term{morphism}{态射}，$W$ 是 $X$ 上\Term{rank}{秩}为 $c$ 的\Term{locally-free}{局部自由}\Term{sheaf}{层}，$F$ 是 $X$ 上的\Term{coherent}{凝聚}\Term{sheaf}{层}，$d$ 是一个 $>0$ 的整数，$s:W\longrightarrow\mathcal O_X$ 是 $W^\vee$ 的一个\Term{section}{截面}。假设 $S$ 是\Term{noetherian}{诺特}的，并且：

(a) 由 $s$ 的消没定义的 $X$ 的\Term{subscheme}{子概形} $H$ 在 $S$ 上\Term{smooth}{光滑}。

(b) 在 $X$ 上\Term{local-scheme}{局部}地，若 $s_i$ 是 $s$ 在 $W^\vee$ 的一组基下的坐标，则 $s_i$ 构成一个 $\mathcal O_X$-\Term{regular-sequence}{正则序列}和 $F$-\Term{regular-sequence}{正则序列}。

(c) 对任意不全为零的整数列 $k_i$，均有
\[
 R^if_*\left(\bigotimes_i\bigwedge^{k_i}W\otimes\Omega^j_{X/S}\otimes F\right)=0\qquad\text{当 }i+j<d.
\]
则有：
\[
\begin{array}{ll}
\text{(i)}&R^if_*(\Omega^j_{H/S}\otimes F)\xrightarrow{\sim}R^if_*(\Omega^j_{X/S}\otimes F),\quad i+j<d-c,\\[3pt]
\text{(ii)}&R^if_*(\Omega^j_{H/S}\otimes F)\lhook\joinrel\longrightarrow R^if_*(\Omega^j_{X/S}\otimes F),\quad i+j=d-c.
\end{array}
\]
\end{SGAresult}
\begin{proof}
$H$ 在 $X$ 中的\Term{conormal-sheaf}{余法层} $N_{H/X}$ \Term{isomorphism}{同构}于 $W$ 在 $H$ 上的\Term{restriction}{限制}，因此\Term{exact-sequence}{正合列}
\[
0\longrightarrow N_{H/X}\longrightarrow\Omega^1_{X/S}\otimes\mathcal O_Y\longrightarrow\Omega^1_{H/S}\longrightarrow0\tag{1.3.1}
\]
在 $\Omega^j_{X/S}\otimes\mathcal O_Y$ 上定义一个递降\Term{filtration}{滤过}，其逐次商为
\[
\operatorname{Gr}^k(\Omega^j_{X/S}\otimes\mathcal O_Y)\simeq\bigwedge^kW\otimes\Omega^{j-k}_{H/S}.\tag{1.3.2}
\]
由 (b) 可知仍有
\[
\operatorname{Gr}^k(\Omega^j_{X/S}\otimes\mathcal O_Y\otimes F)=\bigwedge^kW\otimes\Omega^{j-k}_{H/S}\otimes F.\tag{1.3.3}
\]
我们对 $j$ 作归纳，同时对所有满足 (b)、(c) 的 $F$ 证明 (i)、(ii)。对于 $j<0$，这些断言为空；
\SourcePage{55}
故设对 $j'<j$ 已证得 (i)、(ii)。归纳假设适用于\Term{sheaf}{层} $\bigwedge^kW\otimes F$；结合 (c)，给出
\[
R^if_*\left(\bigwedge^kW\otimes\Omega^{j-k}_{H/S}\otimes F\right)=0\quad\text{当 }k>0\text{ 且 }i+j<d-c+k.\tag{1.3.4}
\]
上述 $\Omega^j_{X/S}\otimes\mathcal O_Y\otimes F$ 的\Term{filtration}{滤过}定义一个\Term{spectral-sequence}{谱序列}
\[
E_1^{p,i-p}=R^if_*\left(\bigwedge^pW\otimes\Omega^{j-p}_{H/S}\otimes F\right)\Longrightarrow R^if_*(\Omega^j_{X/S}\otimes\mathcal O_Y\otimes F).
\]
由 (1.3.4)，当 $p>0$ 且 $i+j\leq d-c$ 时，$E_1^{p,i-p}=0$，从而得到
\[
\left\{\begin{aligned}
R^if_*(\Omega^j_{X/S}\otimes\mathcal O_Y\otimes F)&\xrightarrow{\sim}R^if_*(\Omega^j_{H/S}\otimes F)&&\text{当 }i+j<d-c,\\
R^if_*(\Omega^j_{X/S}\otimes\mathcal O_Y\otimes F)&\lhook\joinrel\longrightarrow R^if_*(\Omega^j_{H/S}\otimes F)&&\text{当 }i+j=d-c.
\end{aligned}\right.\tag{1.3.5}
\]
根据假设 (b)，选取适当的\Term{differential}{微分}后，$\bigwedge^kW\otimes\Omega^j_{X/S}\otimes F$ 构成 $\Omega^j_{X/S}\otimes\mathcal O_Y\otimes F$ 的一个左\Term{koszul-resolution}{Koszul 分解}。这个\Term{filtration}{滤过}定义一个\Term{spectral-sequence}{谱序列}
\[
E_1^{-p,q}=R^qf_*\left(\bigwedge^pW\otimes\Omega^j_{X/S}\otimes F\right)\Longrightarrow R^{q-p}f_*(\Omega^j_{X/S}\otimes\mathcal O_Y\otimes F).\tag{1.3.6}
\]
根据假设 (c)，当 $p>0$ 且 $q+j<d$ 时，$E_1^{-p,q}=0$；当 $p\notin[0,c]$ 时，它也显然为零。因此，
\[
\left\{\begin{aligned}
R^if_*(\Omega^j_{X/S}\otimes F)&\xrightarrow{\sim}R^if_*(\Omega^j_{H/S}\otimes F)&&\text{当 }i+j<d-c,\\
R^if_*(\Omega^j_{X/S}\otimes F)&\lhook\joinrel\longrightarrow R^if_*(\Omega^j_{H/S}\otimes F)&&\text{当 }i+j=d-c.
\end{aligned}\right.\tag{1.3.7}
\]
现在由 (1.3.5)、(1.3.7) 即得断言 (i)、(ii)。
\end{proof}

\SourcePage{56}
\textbf{1.4。}设 $E$ 是\Term{scheme}{概形} $S$ 上\Term{rank}{秩}为 $r+1$ 的\Term{locally-free}{局部自由}模\Term{sheaf}{层}，$\mathbf P(E)$ 为相应的\Term{projective-bundle}{射影丛}，$d$ 为一个 $\leq r$ 的整数，$\underline a=(a_1,\ldots,a_d)$ 为一个整数列。$\mathbf P(E)$ 中一个\Term{multidegree}{多重次数}为 $\underline a$ 的\Term{relative-complete-intersection}{相对完全交}，是指 $\mathbf P(E)$ 的一个\Term{subscheme}{子概形}，它在每个\Term{fibre}{纤维}中都有\Term{codimension}{余维} $d$，且在 $S$ 的\Term{etale-topology}{平展拓扑}下\Term{local-scheme}{局部}地，是 $d$ 个\Term{degree}{次数}分别为 $a_1,\ldots,a_d$ 的\Term{hypersurface}{超曲面}的交。

由 1.1，取 $X/S=\mathbf P(E)/S$、$F=\mathcal O_X$、$W=\bigoplus_i\mathcal O(-a_i)$，并取 $s$ 为 $W^\vee$ 的一个定义\Term{multidegree}{多重次数}为 $\underline a$ 的\Term{relative-complete-intersection}{相对完全交} $Y\subset\mathbf P(E)$ 的\Term{section}{截面}时，1.3 的假设 (b)、(c) 均满足。

\begin{SGAresult}{\Term{theorem}{定理}}{1.5}
在 1.4 的假设下，设 $n=r-d$ 是 $Y$ 的\Term{relative-dimension}{相对维数}，$\eta\in H^0(S,R^1f_*\Omega^1_{Y/S})$ 是 $\mathcal O(1)$ 的\Term{first-chern-class}{第一 Chern 类}，其中
\[
\begin{tikzcd}
Y\arrow[rr,hook]\arrow[dr,"f"']&&\mathbf P(E)\arrow[dl,"g"]\\
&S&
\end{tikzcd}
\]
若 $Y$ 在 $S$ 上\Term{smooth}{光滑}，则：

(i) $R^if_*\Omega^j_{Y/S}$ 是\Term{locally-free}{局部自由}的，其形成与任意\Term{base-change}{基变换}相容。

(ii) 联系\Term{hodge-cohomology}{Hodge 上同调}和\Term{de-rham-cohomology}{De Rham 上同调}的\Term{spectral-sequence}{谱序列}
\[
 E_1^{pq}=R^qf_*\Omega^p_{Y/S}\Longrightarrow R^{p+q}f_*(\Omega^\bullet_{Y/S})
\]
\Term{degeneration}{退化}（$E_1=E_\infty$），其\Term{abutment}{终项}是\Term{locally-free}{局部自由}的，且其形成与任意\Term{base-change}{基变换}相容。

$(\mathrm{iii})_a$ 有
\[
 R^if_*\Omega^j_{Y/S}=0\qquad\text{若 }i\ne j\text{ 且 }i+j\ne n.
\]
\SourcePage{57}
$(\mathrm{iii})_b$ 若 $2i<n$，则 $\eta^i$ 是 $R^if_*\Omega^i_{Y/S}$ 的一组基。

$(\mathrm{iii})_c$ 若 $n<2i\leq2n$，且 $\alpha_i$ 是 $R^if_*\Omega^i_{Y/S}$ 中与 $R^{n-i}f_*\Omega^{n-i}_{Y/S}$ 的基 $\eta^{n-i}$ \Term{dual-basis}{对偶的基}，则
\[
 \eta^i=\left(\prod_{1}^{d}a_i\right)\alpha_i.
\]
$(\mathrm{iii})_d$ 若 $n=2m$，则 $\eta^m\ne0$（当 $S\ne\varnothing$ 时），并且商 $R^mf_*\Omega^m_{Y/S}/\mathcal O_S\cdot\eta^m$ 是\Term{locally-free}{局部自由}的。
\end{SGAresult}
\begin{proof}
只需在一个“\Term{universal}{普遍}”情形下证明 1.5，因此可假设 $S$ 在 $\operatorname{Spec}(\mathbf Z)$ 上\Term{smooth}{光滑}。此时，断言 (ii) 由 (i) 得到，因为根据\Term{hodge-theory}{Hodge 理论}，它在 $\mathbf C$ 上成立，从而在任意\Term{characteristic-field}{特征}为零的域上成立。此外，由于 $S$ 是\Term{reduced}{约化}的，假设对 $S$ 为域的\Term{spectrum}{谱}时已证得 (iii)，则 (i) 由 (iii) 得到：事实上，由 (iii)，对每个 $j$，函数
\[
 s\longmapsto\dim_{k(s)}H^i(Y_s,\Omega^j_{Y_s})
\]
在 $S$ 上恒定，至多 $i=j$ 时例外；由于它们的\Term{alternating-sum}{交错和}恒定（EGA III 7.9.2），它们便全都恒定，再应用 EGA III 6.10.5 + 7.6.9 + 7.7.5 即得结论。

由 (1.4)、(1.3)，有
\[
\left\{\begin{aligned}
R^if_*\Omega^j_{Y/S}&\xrightarrow{\sim}R^ig_*\Omega^j_{\mathbf P(E)/S}&&\text{当 }i+j<n,\\
R^if_*\Omega^j_{Y/S}&\lhook\joinrel\longrightarrow R^ig_*\Omega^j_{\mathbf P(E)/S}&&\text{当 }i+j=n.
\end{aligned}\right.\tag{1.5.1}
\]
若 $S$ 是域的\Term{spectrum}{谱}，则断言 $(\mathrm{iii})_a$、$(\mathrm{iii})_b$、$(\mathrm{iii})_d$ 由 (1.5.1) 和\Term{serre-duality}{Serre 对偶性}（[2] V 11.2，第 211 页）得到。
\SourcePage{58}
\Term{cohomology-class}{上同调类} $\eta^n$ 是一个\Term{degree}{次数}为 $\prod_1^d a_i$（即\Term{complete-intersection}{完全交}的\Term{degree}{次数}）的 $0$-\Term{cycle}{圈}的\Term{cohomology-class}{上同调类}，从而
\[
 \eta^n=\left(\prod_1^d a_i\right)\alpha_n.
\]
由定义，$\eta^{n-i}\alpha_i=\alpha_n$，由此即得 $(\mathrm{iii})_c$。

因此，断言 (i)、(ii) 成立；而在\Term{universal}{普遍}情形下，断言 (iii) 由 (i) 以及假设对 $S$ 为域的\Term{spectrum}{谱}时已证得的 (iii) \Term{pushout}{推出}。
\end{proof}

\begin{SGAresult}{\Term{theorem}{定理}}{1.6}
保留 1.4、1.5 的假设和记号，假设 $Y$ \Term{smooth}{光滑}，$\ell$ 是与 $S$ 的所有\Term{residue-characteristic}{剩余特征}互素的素数，以
\[
 \eta_\ell\in H^0(S,R^2f_*(\mathbf Z_\ell(1)))
\]
表示 $\mathcal O(1)$ 的 $\ell$ 进\Term{first-chern-class}{第一 Chern 类}（SGA 5 VII）。则：

(i) 若 $i$ 为奇数且不等于 $n$，则 $R^if_*\mathbf Z_\ell=0$。

(ii) 若 $2i<n$，则 $R^{2i}f_*\mathbf Z_\ell(i)$ \Term{isomorphism}{同构}于 $\mathbf Z_\ell$，且以 $\eta_\ell^i$ 为基。

(iii) 若 $n<2i\leq2n$，则 $R^{2i}f_*\mathbf Z_\ell(i)$ \Term{isomorphism}{同构}于 $\mathbf Z_\ell$，且以 $\eta_\ell^i/\prod_1^d a_i$ 为基。

(iv) $\mathbf Z_\ell$-\Term{sheaf}{层} $R^nf_*\mathbf Z_\ell$ 是\Term{locally-constant}{局部常值}、\Term{constructible}{可构造}且\Term{torsion-free}{无挠}的；若 $n$ 为偶数，则 $\eta_\ell^{n/2}\ne0$（当 $S\ne\varnothing$ 时），并且 $R^nf_*\mathbf Z_\ell(n/2)/\mathbf Z_\ell\cdot\eta_\ell^{n/2}$ \Term{torsion-free}{无挠}。
\end{SGAresult}
\begin{proof}
根据 SGA 4 XVI 2.2，$\mathbf Z_\ell$-\Term{sheaf}{层} $R^if_*\mathbf Z_\ell$ 是\Term{locally-constant}{局部常值}\Term{constructible}{可构造}的，其形成与任意\Term{base-change}{基变换}相容。因此，只需在 $S$ 为\Term{algebraically-closed}{代数闭}域的\Term{spectrum}{谱}时证明 1.6。反复应用 SGA 4 XIV 3.3，得到对任意 $m$，\Term{canonical}{典范}\Term{arrow}{箭头}
\[
 H_Y^q(\mathbf P(E),\mathbf Z/\ell^m)\longrightarrow H^q(Y,\mathbf Z/\ell^m)
\]
\SourcePage{59}
在 $q>r+d$ 时为\Term{bijection}{双射}，在 $q=r+d$ 时为\Term{surjectivity}{满射}。由此和\Term{poincare-duality}{Poincaré 对偶性}（SGA 4 XVIII），可\Term{pushout}{推出}（SGA 5 VII）\Term{restriction}{限制}\Term{arrow}{箭头}
\[
 H^q(\mathbf P(E),\mathbf Z_\ell)\longrightarrow H^q(Y,\mathbf Z_\ell)
\]
在 $q<n$ 时为\Term{bijection}{双射}，在 $q=n$ 时为\Term{injectivity}{单射}且\Term{cokernel}{余核}\Term{torsion-free}{无挠}。这就证明了 $i<n$ 时的 (i)，以及 (ii)、(iv)。与 1.5 一样可验证，$\eta_\ell^n$ 是\Term{canonical}{典范}类的 $\prod_1^d a_i$ 倍，于是再由\Term{poincare-duality}{Poincaré 对偶性}（SGA 4 XVIII）得到 (i)、(iii)。
\end{proof}

\textbf{1.7。}若 $Y$ 是域 $k$ 上的\Term{projective-space}{射影空间} $\mathbf P^r(k)$ 中\Term{dimension}{维数}至少为 $3$ 的\Term{complete-intersection}{完全交}，则已知 $\operatorname{Pic}(Y)$ \Term{isomorphism}{同构}于 $\mathbf Z$，并由 $\mathcal O(1)$ 生成（SGA 2 XII 3.7）。当\Term{dimension}{维数}为 $2$ 且 $Y$ \Term{smooth}{光滑}时，还能证明：
\begin{SGAresult}{\Term{theorem}{定理}}{1.8}
设 $Y$ 是\Term{algebraically-closed}{代数闭}域 $k$ 上的\Term{projective-space}{射影空间} $\mathbf P^r(k)$ 中的\Term{smooth}{光滑}\Term{complete-intersection}{完全交}\Term{surface}{曲面}。则 $\operatorname{Pic}^\tau(Y)=0$，且 $\mathcal O(1)$ 的类在\Term{neron-severi-group}{Néron--Severi 群} $\operatorname{NS}(Y)$ 中不可整除。
\end{SGAresult}
\begin{proof}
由于 $H^1(Y,\mathcal O)=0$（1.5 $(\mathrm{iii})_a$），已知 $\operatorname{Pic}^0(Y)=0$。于是，由下面的著名\Term{lemma}{引理}、$H^2(Y,\mathbf Z_\ell(1))$ \Term{torsion-free}{无挠}这一事实（1.6 (iv)），以及 $H^0(Y,\Omega_Y^1)=0$ 这一事实（1.5 $(\mathrm{iii})_a$），即可得 $\operatorname{Pic}^\tau(Y)=0$。
\end{proof}
\begin{SGAresult}{\Term{lemma}{引理}}{1.9}
设 $Y$ 是\Term{algebraically-closed}{代数闭}域 $k$ 上的\Term{complete-variety}{完备代数簇}。

(i) 若 $\ell$ 与 $k$ 的\Term{characteristic-field}{特征}互素，则 $\operatorname{NS}(Y)$ 的 $\ell$-\Term{torsion}{挠子群}可与 $H^2(Y,\mathbf Z_\ell(1))$ 的 $\ell$-\Term{torsion}{挠子群}相等同。
\SourcePage{60}
(ii) 若 $k$ 的\Term{characteristic-field}{特征}为非零的 $p$，且 $Y$ \Term{smooth}{光滑}，则 $\operatorname{Pic}(Y)$ 中乘以 $p$ 的\Term{kernel}{核}可与 $H^0(Y,\Omega_Y^1)$ 中由\Term{locally-logarithmic-differential}{局部对数微分}构成的\Term{subgroup}{子群}相等同。
\end{SGAresult}
\begin{proof}
\Term{kummer-sequence}{Kummer 正合列}
\[
0\longrightarrow\mu_{\ell^n}\longrightarrow\mathbf G_m\xrightarrow{x\mapsto x^{\ell^n}}\mathbf G_m\longrightarrow0
\]
定义\Term{exact-sequence}{正合列}
\[
0\longrightarrow\operatorname{Pic}(Y)/\ell^n\operatorname{Pic}(Y)\longrightarrow H^2(Y,\mu_{\ell^n})\longrightarrow(H^2(\mathbf G_m))_{\ell^n}\longrightarrow0,
\]
其中 $\operatorname{Pic}(Y)/\ell^n\operatorname{Pic}(Y)=\operatorname{NS}(Y)\otimes\mathbf Z/\ell^n$。取\Term{projective-limit}{投射极限}，得到
\[
0\longrightarrow\operatorname{NS}(Y)\otimes\mathbf Z_\ell\longrightarrow H^2(Y,\mathbf Z_\ell(1))\longrightarrow\operatorname{Hom}(\mathbf Q_\ell/\mathbf Z_\ell,H^2(\mathbf G_m))\longrightarrow0,
\]
由于 $\mathbf Q_\ell/\mathbf Z_\ell$ 是\Term{divisible-group}{可除的}，\Term{cokernel}{余核}\Term{torsion-free}{无挠}。断言 (i) 随之得到。

为证明 (ii)，只需注意\Term{exact-sequence}{正合列}
\[
0\longrightarrow\mathcal O_Y^*\xrightarrow{x\mapsto x^p}\mathcal O_Y^*\xrightarrow{df/f}\Omega_Y^1
\]
定义一个\Term{short-exact-sequence}{短正合列}
\[
0\longrightarrow\mathcal O_Y^*\longrightarrow\mathcal O_Y^*\longrightarrow\Omega^1_{Y,\mathrm{loc\ log}}\longrightarrow0,
\]
然后取相应的\Term{long-exact-sequence}{上同调长正合列}。这个论证归功于 Cartier；他还证明，若以 $C$ 表示 $H^0(Y,\Omega_Y^1)$ 上的\Term{cartier-operation}{Cartier 运算}，则\Term{locally-logarithmic-differential}{局部对数微分}构成 $C-\operatorname{Id}$ 的\Term{kernel}{核}。
\end{proof}
\begin{proof}[\Term{theorem}{定理} 1.8 的证明（续）]
反设 $\mathcal O(1)$ 在 $\operatorname{NS}(Y)$ 中的类形如 $nu$，其中 $n>1$。若 $\ell$ 整除 $n$ 且与\Term{characteristic-field}{特征}互素，
\SourcePage{61}
则在 $\ell$ 进\Term{cohomology}{上同调}中仍有
\[
 \eta_\ell=n\,c_1(u),
\]
这与 1.6 (iv) 矛盾。

最后，若 $k$ 的\Term{characteristic-field}{特征} $p$ 非零且 $p$ 整除 $n$，则在\Term{hodge-cohomology}{Hodge 上同调}中仍有
\[
 \eta=n\,c_1(u)=0,
\]
这与 1.5 $(\mathrm{iii})_d$ 矛盾，证明完成。
\end{proof}

\SGAsection{2}{数值结果}
\textbf{2.1。}设 $\underline a=(a_i)_{1\leq i\leq d}$ 是一个 $\geq1$ 的整数列，以 $V_n(\underline a)$ 表示域 $k$ 上\Term{projective-space}{射影空间} $\mathbf P^{n+d}$ 中\Term{dimension}{维数}为 $n$、\Term{multidegree}{多重次数}为 $\underline a$ 的一个\Term{smooth}{光滑}\Term{complete-intersection}{完全交}。令
\[
 h^{pq}(\underline a)=\dim_k H^q(V_{p+q}(\underline a),\Omega^p)
\]
及
\[
 h_o^{pq}(\underline a)=h^{pq}(\underline a)-\delta_{p,q}.
\]
于是，$h_o^{pq}(\underline a)$ 是\Term{multidegree}{多重次数}为 $\underline a$ 的\Term{complete-intersection}{完全交}的中维\Term{primitive-cohomology}{本原上同调}的\Term{hodge-number}{Hodge 数}。它们只依赖于 $\underline a,p,q$（1.5）。令
\[
 H(\underline a)=\sum_{p,q\geq0}h_o^{p,q}(\underline a)y^pz^q\in\mathbf Z[[y,z]].\tag{2.1.1}
\]
当 $d=1$ 且 $\underline a=(a)$ 时，令 $h^{pq}(a)=h^{pq}(\underline a)$、$H(a)=H(\underline a)$。当 $d=0$ 时，令 $\sup(a_i)=1$。

下面的公式是 Hirzebruch [3] 第 160 页公式 (2) 在 $k=0$ 时的特殊情形：
\SourcePage{62}
\[
\sum_{n=0}^{\infty}\sum_p\chi(V_n(\underline a),\Omega^p)y^pz^{n+d}
=\frac{1}{(1+yz)(1-z)}\prod_{1}^{d}\frac{(1+yz)^{a_i}-(1-z)^{a_i}}{(1+yz)^{a_i}+y(1-z)^{a_i}}.\tag{2.2}
\]
我们要把这个公式变换为下式：
\begin{SGAresult}{\Term{theorem}{定理}}{2.3（Hirzebruch）}
使用 2.1 的记号，有
\[
H(\underline a)=\frac{1}{(1+y)(1+z)}\left[\prod_{1\leq i\leq d}\frac{(1+y)^{a_i}-(1+z)^{a_i}}{-(1+y)^{a_i}z+(1+z)^{a_i}y}-1\right].
\]
\end{SGAresult}
\begin{proof}
由 1.5，有
\[
\chi(V_n(\underline a),\Omega^p)=(-1)^p+(-1)^{n-p}h_o^{p,n-p}(\underline a),
\]
所以 (2.2) 的左端 $\mathrm I$ 可改写为
\[
\begin{aligned}
\mathrm I
&=\sum_{n=0}^{\infty}\sum_{p+q=n}\bigl((-1)^p+(-1)^qh_o^{p,q}(\underline a)\bigr)y^pz^{p+q+d}\\
&=\sum_{p,q\geq0}(-1)^p(yz)^pz^qz^d+\sum_{p,q\geq0}(-1)^qh_o^{p,q}(\underline a)(yz)^pz^qz^d\\
&=z^d\left[\frac{1}{(1+yz)(1-z)}+H(\underline a)(yz,-z)\right].
\end{aligned}
\]
(2.2) 的右端 $\mathrm {II}$ 可改写为
\[
\mathrm {II}=\frac{z^d}{(1+yz)(1-z)}\prod_1^d\frac{(1+yz)^{a_i}-(1-z)^{a_i}}{(1+yz)^{a_i}z+yz(1-z)^{a_i}},
\]
于是，在 $\mathrm I=\mathrm {II}$ 中约去 $z^d$，得到
\[
H(\underline a)(yz,-z)=\frac{1}{(1+yz)(1-z)}\left[\prod_1^d\frac{(1+yz)^{a_i}-(1-z)^{a_i}}{(1+yz)^{a_i}z+(1-z)^{a_i}yz}-1\right],
\]
再作变量替换，即由此式得到 2.3。
\end{proof}

\SourcePage{63}
\begin{SGAresult}{推论}{2.4}
使用 2.1 的记号，

(i) 当 $d=1$ 时，有
\[
\begin{aligned}
H(a)&=-\frac{(1+y)^{a-1}-(1+z)^{a-1}}{(1+y)^az-(1+z)^ay}\\
&=\frac{\displaystyle\sum_{i,j\geq0}\binom{a-1}{i+j+1}y^iz^j}{\displaystyle1-\sum_{i,j\geq1}\binom{a}{i+j}y^iz^j}.
\end{aligned}
\]
(ii) 以 $|P|$ 表示 $[1,d]$ 的\Term{subset}{子集} $P$ 的元素个数，则有
\[
H(\underline a)=\sum_{\substack{P\subset[1,d]\\P\ne\varnothing}}(1+y)^{|P|-1}(1+z)^{|P|-1}\prod_{i\in P}H(a_i).
\]
\end{SGAresult}
\begin{proof}
根据 (2.3)，有
\[
\begin{aligned}
H(a)&=\frac{-1}{(1+y)(1+z)}\left[\frac{(1+y)^a-(1+z)^a}{(1+z)^a\cdot z-(1+z)^ay}+1\right]\\
&=\frac{-1}{(1+y)(1+z)}\left[\frac{(1+y)^a(1+z)-(1+z)^a(1+y)}{(1+y)^az-(1+z)^ay}\right]\\
&=-\frac{(1+y)^{a-1}-(1+z)^{a-1}}{(1+y)^az-(1+z)^ay}.
\end{aligned}
\]
以 $N,D$ 分别表示此式的分子和分母。有
\[
 H(a)=-\frac ND,
\]
\SourcePage{64}
\[
\begin{aligned}
N&=\sum_k\binom{a-1}{k}(y^k-z^k)
=(y-z)\sum_{i,j\geq0}\binom{a-1}{i+j+1}y^iz^j,\\
D&=\sum_k\binom ak(y^kz-yz^k)
=(z-y)+yz\sum_{k\geq1}\binom ak(y^{k-1}-z^{k-1})\\
&=(y-z)\left(-1+yz\sum_{i,j\geq0}\binom a{i+j+2}y^iz^j\right)\\
&=(y-z)\left(-1+\sum_{i,j\geq1}\binom a{i+j}y^iz^j\right),
\end{aligned}
\]
断言 (i) 随之得到。

当 $d=1$ 时，反解 2.3 的公式，得到
\[
-\frac{(1+y)^a-(1+z)^a}{(1+y)^az-(1+z)^ay}=1+(1+x)(1+y)H(a).
\]
将此式代入 2.3，得到
\[
\begin{aligned}
H(\underline a)
&=\frac1{(1+y)(1+z)}\left[\prod_1^d\bigl(1+(1+x)(1+y)H(a_i)\bigr)-1\right]\\
&=\frac1{(1+y)(1+z)}\left[\sum_{\substack{P\subset[1,d]\\P\ne\varnothing}}(1+x)^{|P|}(1+y)^{|P|}\prod_{i\in P}H(a_i)\right],
\end{aligned}
\]
由此得 (ii)。
\end{proof}

\begin{SGAresult}{\Term{theorem}{定理}}{2.5}
使用 2.1 的记号：

(i) 若 $p+q=p'+q'$ 且 $p\leq p'\leq q'\leq q$，则
\[
 h_o^{pq}(\underline a)\leq h_o^{p'q'}(\underline a).
\]
(ii) 若 $p\leq q$，则 $h_o^{pq}(\underline a)\ne0$ 的充要条件是
\SourcePage{65}
\[
 p>\left[\frac{p+q+d-\sum a_i}{\sup(a_i)}\right],
\]
其中 $[\ ]$ 表示商的整数部分。
\end{SGAresult}
\begin{proof}
以 (*) 表示\Term{formal-power-series}{形式幂级数} $P=\sum a_{pq}y^pz^q$ 的下述性质：

(*) $a_{pq}$ 均非负，$a_{pq}=a_{qp}$，且若 $p+q=p'+q'$、$p\leq p'\leq q'\leq q$，则 $a_{pq}\leq a_{p'q'}$。

断言 (i) 立即由公式 (2.4)、恒等式 $(1-x)^{-1}=\sum_{n\geq0}x^n$ 以及下面的\Term{lemma}{引理}得到，\Term{lemma}{引理}的证明留给读者：
\end{proof}
\begin{SGAresult}{\Term{lemma}{引理}}{2.5.1}
满足 (*) 的\Term{formal-power-series}{形式幂级数}的\Term{set}{集合}对加法和\Term{multiplication}{乘法}封闭。
\end{SGAresult}
\begin{proof}[\Term{theorem}{定理} 2.5 的证明（续）]
首先在\Term{hypersurface}{超曲面}（$d=1$）这一特殊情形下证明 (ii)。此时，对于 $p\leq q$，待证公式可改写为
\[
 h_o^{pq}(a)\ne0\quad\Longleftrightarrow\quad p\geq\left[\frac{p+q+1}{a}\right]\tag{2.5.2}
\]
（这又显然\Term{equivalence}{等价}于 $(p+1)a\geq p+q+2$）。

由 2.4 (i) 的第二个公式，得到 $h_o^{pq}$ 非零，当且仅当存在数对 $(e,f)$ 和 $(c_i,d_i)$，使得
\[
\left\{\begin{aligned}
{}&(e,f)\geq(0,0)\quad\text{且 }e+f\leq a-2,\\
&(c_i,d_i)\geq(1,1)\quad\text{且 }c_i+d_i\leq a,\\
&(p,q)=(e,f)+\sum(c_i,d_i).
\end{aligned}\right.\tag{2.5.3}
\]
\SourcePage{66}
给定 $p+q$，最小的可能的 $p$ 在 $e=0$、$c_i=1$ 时取得，且等于指标 $i$ 的个数。因此，这个 $p$ 是满足
\[
 p\geq\frac{(p+q)-(a-2)}a,\quad\text{即 }(p+1)a\geq p+q+2
\]
的所有 $p$ 的\Term{infimum}{下确界}。

转到\Term{generic}{一般}情形。若 $\underline a'$ 是从 $\underline a$ 中删去所有等于 $1$ 的 $a_i$ 而得，则 $H(\underline a)=H(\underline a')$。由公式 2.5 (ii) 的\Term{form}{形式}，可假设对每个 $i$ 都有 $a_i>1$。若 $d=0$，则 $h_o^{pq}=0$，断言成立；故可不失\Term{generic}{一般}性地假设 $d>0$ 且 $a_d=\sup(a_i)$。

由 2.4 (ii)、(2.5.2)，$h_o^{pq}(\underline a)\ne0$ 的充要条件是存在 $[1,d]$ 的一个\Term{nonempty}{非空}\Term{subset}{子集} $P$、一个数对 $(e,f)$ 和数对 $(c_i,d_i)_{1\leq i\leq d}$，使得
\[
\left\{\begin{aligned}
{}&(|P|-1,|P|-1)\geq(e,f)\geq(0,0),\\
&c_i,d_i\geq\left[\frac{c_i+d_i+1}{a_i}\right];\quad(c_i,d_i)\geq(0,0),\\
&(p,q)=(e,f)+\sum(c_i,d_i).
\end{aligned}\right.\tag{2.5.4}
\]
固定 $n=p+q$，来求 $p$ 的最小可能值。若 $n_i=c_i+d_i$，则对于这样的 $p$，条件 (2.5.4) 用 $c_i,n_i,e,f$ 可改写为
\SourcePage{67}
\[
\left\{\begin{aligned}
{}&a_i=2\ \Longrightarrow\ n_i\text{ 为偶数}\quad\text{（否则 (2.5.4) 不可能成立）},\\
&c_i\geq\left[\frac{n_i+1}{a_i}\right];\quad c_i\leq n_i,\\
&(|P|-1,|P|-1)\geq(e,f)\geq(0,0),\\
&n=e+f+\sum n_i,\\
&p=e+\sum c_i.
\end{aligned}\right.\tag{2.5.5}
\]
这个描述首先表明，最小的 $p$ 在 $P=[1,d]$ 时取得，以下仅考虑这种情形。若
\[
 \sum(a_i-2)+(d-1)\geq n,\tag{2.5.6}
\]
则存在一个满足 $p=c_i=e=0$ 的解。否则，若 $a_d\ne2$，最佳选择是对 $i\ne d$ 取 $n_i=a_i-2$，从而对 $i\ne d$ 有 $c_i=0$，并且 $e=0$、$f=d-1$、$n_d=n-(d-1)-\sum_{i\ne d}(a_i-2)$。于是
\[
 p=\left[\frac{n-(d-1)-\sum_{i\ne d}(a_i-2)+1}{a_d}\right],\tag{2.5.7}
\]
所以，对于 $p\leq q$，有 $h_o^{pq}(\underline a)\ne0$，当且仅当
\[
 p\geq\left[\frac{n-\sum_{i\ne d}(a_i-1)+1}{a_d}\right].\tag{2.5.8}
\]
当 $d=1$ 时，公式 (2.5.8) 化为 (2.5.2)；当 (2.5.6) 成立时，它显然成立。若 $d\geq2$、(2.5.6) 不成立，且所有 $a_i$ 均等于 $2$，最佳选择仍然是
\SourcePage{68}
取 $n_i=0$（$i\ne d$）、$e=0$，$f$ 取 $d-1$ 或 $d-2$ 中模 $2$ 与 $n$ 同余者，并取 $n_d=n-f$。若 $n\equiv d-1\pmod2$，可像上面一样证明 (2.5.7)。否则，注意到
\[
 \left[\frac{n-(d-2)+1}{2}\right]=\left[\frac{n-(d-1)+1}{2}\right],
\]
故即使 $a_d=2$，(2.5.7)，从而 (2.5.8)，也仍成立。

剩下只需验证 (2.5.8) \Term{equivalence}{等价}于公式 2.5 (ii)。事实上，有
\[
\begin{aligned}
\left[\frac{n-\sum_{i\ne d}(a_i-1)+1}{a_d}\right]
&=\left[\frac{n+d-\sum_{i\ne d}a_i}{a_d}\right]\\
&=\left[\frac{p+q+d-\sum_i a_i}{a_d}\right]+1.
\end{aligned}
\]
\end{proof}

\textbf{2.6。}若 $X$ 是 $\mathbf P^r$ 中的一个 $n$ 维\Term{complete-intersection}{完全交}，则由 (2.5.8) 立即得到：$H^n(X)$ 的\Term{primitive-cohomology}{本原部分}只有在 $X$ 是奇数维\Term{quadric}{二次超曲面}时才为零。

\begin{SGAresult}{定义}{2.7}
若 $X$ 是 $\mathbf C$ 上的\Term{proper}{固有}\Term{smooth}{光滑}\Term{algebraic-variety}{代数簇}，则 $H^n(X)$ 的\Term{hodge-level}{Hodge 水平}是所有满足 $p+q=n$ 且 $h^{pq}(X)\ne0$ 的数 $q-p$ 的\Term{supremum}{上确界}。
\end{SGAresult}
因此，\Term{hodge-level}{Hodge 水平} $\ell$ 模 $2$ 与 $n$ 同余，且至多等于 $\inf(n,\dim(X)-n)$。除 $H^n(X)=0$ 时等于 $-\infty$ 外，它总是 $\geq0$。

\begin{SGAresult}{推论}{2.8}
采用 2.1 的记号，取 $k=\mathbf C$，并假设 $a_d=\sup(a_i)$。设 $\ell$ 是满足 $0\leq\ell<n$ 且模 $2$ 与 $n$ 同余的整数。$H^n(V_n(\underline a))$ 的\Term{hodge-level}{Hodge 水平} $\leq\ell$ 的充要条件是
\SourcePage{69}
\[
 \sum_{i\ne d}(a_i-1)\leq(\ell+1)-(a_d-2)\frac{n-\ell}{2}.
\]
\end{SGAresult}
\begin{proof}
由 (2.5.8)，若\Term{hodge-level}{Hodge 水平} $\geq0$，它便是 $0$ 与
\[
 n-2\left[\frac{n-\sum_{i\ne d}(a_i-1)+1}{a_d}\right]
\]
的\Term{supremum}{上确界}，而不等式
\[
 n-2\left[\frac{n-\sum_{i\ne d}(a_i-1)+1}{a_d}\right]\leq\ell
\]
可改写为
\[
 \frac{n-\ell}{2}\leq\left[\frac{n-\sum_{i\ne d}(a_i-1)+1}{a_d}\right],
\]
或改写为
\[
 \sum_{i\ne d}(a_i-1)\leq1+n-a_d\cdot\frac{n-\ell}{2}
 =\ell+1-(a_d-2)\frac{n-\ell}{2}.
\]
\end{proof}

\textbf{2.9。}借助公式 2.8，在 $\ell$ 较小时，很容易列举 $\mathbf P^r$ 中\Term{dimension}{维数} $n>1$、且 $H^n$ 的\Term{hodge-level}{Hodge 水平}为 $\ell$ 的\Term{complete-intersection}{完全交}的\Term{multidegree}{多重次数}。不过，本讲初稿中的表格有误；下面的表格摘自 M. Rapoport 在其文章《对 P. Deligne〈K3 \Term{surface}{曲面}的 Weil 猜想〉一文的补充》（\textit{Complément à l'article de P. Deligne ``La conjecture de Weil pour les surfaces K3''}，Inventiones Math. \textbf{15}，227--236（1972））中给出的更完整的表格。对于 $\ell\leq1$，\Term{dimension}{维数}为 $n$、\Term{multidegree}{多重次数}为 $\underline a=(a_1,\ldots,a_d)$，满足 $n>\ell$、$d\geq1$、$a_i\geq2$，且\Term{hodge-level}{Hodge 水平}为 $\ell$ 的\Term{complete-intersection}{完全交} $V_n(\underline a)$ 如下：
\SourcePage{70}
\[
\begin{array}{c@{\qquad}l}
\ell=-\infty&V_{2k+1}(2)\quad(k\geq0),\\[3pt]
\ell=0&V_{2k}(2)\ (k\geq1),\quad V_{2k}(2,2)\ (k\geq1),\quad V_2(3),\\[3pt]
\ell=1&V_{2k+1}(2,2)\ (k\geq1),\quad V_{2k+1}(2,2,2)\ (k\geq1),\\
&V_3(3),\quad V_5(3),\quad V_3(3,2),\quad V_3(4).
\end{array}
\]
\textbf{2.10。}对于较小的 $\underline a$，由 2.4，$h_o^{pq}(\underline a)$ 的\Term{generating-function}{生成函数}可写为
\[
\begin{aligned}
 H(2)&=\frac1{1-yz},\\
 H(3)&=\frac{2+y+z}{1-3yz-y^2z-yz^2},\\
 H(2,2)&=\frac{y+z}{(1-yz)^2}+\frac2{(1-yz)^2}+\frac1{1-yz},
\end{aligned}
\]
因而
\[
\begin{aligned}
 h_o^{pp}(2)&=1,\\
 h_o^{11}(3)&=6,\\
 h_o^{pp}(2,2)&=2p+3.
\end{aligned}
\]

\SourcePage{71}
\SGAsection{}{参考文献}
\begin{enumerate}
\item[{[1]}] Y. Akizuki 与 S. Nakano，关于 Kodaira--Spencer 对 Lefschetz \Term{theorem}{定理}的证明的注记（\textit{Note on Kodaira--Spencer's proof of Lefschetz theorems}），Proc. Jap. Acad. \textbf{30}，266--272（1954）。
\item[{[2]}] R. Hartshorne，\Term{residue}{留数}与\Term{duality}{对偶性}（\textit{Residues and duality}），Lecture Notes \textbf{20}，Springer，1966。
\item[{[3]}] F. Hirzebruch，\Term{algebraic-geometry}{代数几何}中的\Term{topology}{拓扑}方法（\textit{Topological methods in algebraic geometry}），第 3 版，Springer-Verlag，1966。
\item[{[4]}] K. Kodaira，解析叠理论中的一种\Term{differential}{微分}几何方法（\textit{On a differential geometric method in the theory of analytic stacks}），Proc. Nat. Acad. Sci. USA \textbf{39}，1268--1273（1953）。
\end{enumerate}
